\textbf{Phần 1. Trắc nghiệm lựa chọn}
\setcounterpageref{ex}{0}

\begin{ex} %[4D1N01]
Số \textit{Tám mươi lăm nghìn ba trăm linh hai} được viết là
\choice
{$85\,320$}
{$85\,032$}
{$85\,322$}
{\True $85\,302$}
\loigiai{
Số đó viết là $85\,302$.\\
\textbf{Đáp án của câu hỏi là D}
}
\end{ex}

\begin{ex} %[4D1N01]
Trong số $74\,528$, chữ số $4$ có giá trị là
\choice
{\True $4000$}
{$400$}
{$40\,000$}
{$40$}
\loigiai{
Chữ số 4 ở hàng nghìn nên có giá trị là $4000$.\\
\textbf{Đáp án của câu hỏi là A}
}
\end{ex}

\begin{ex} %[4D1H01]
Số liền sau của số $99\,999$ là
\choice
{$99\,998$}
{$100\,001$}
{\True $100\,000$}
{$10\,000$}
\loigiai{
Số liền sau của $99\,999$ là $99\,999 + 1 = 100\,000$.\\
\textbf{Đáp án của câu hỏi là C}
}
\end{ex}

\begin{ex} %[4D1H01]
Số lớn nhất có 5 chữ số khác nhau là
\choice
{$99\,999$}
{$98\,756$}
{$97\,865$}
{\True $98\,765$}
\loigiai{
Số lớn nhất có 5 chữ số khác nhau là $98\,765$.\\
\textbf{Đáp án của câu hỏi là D}
}
\end{ex}

\begin{ex} %[4D1H01]
Số gồm $6$ chục nghìn, $2$ trăm và $5$ đơn vị viết là
\choice
{$62\,005$}
{$60\,250$}
{\True $60\,205$}
{$62\,500$}
\loigiai{
Số gồm $6$ chục nghìn, $0$ nghìn, $2$ trăm, $0$ chục, $5$ đơn vị là $60\,205$.\\
\textbf{Đáp án của câu hỏi là C}
}
\end{ex}

\begin{ex} %[4D1H01]
Sắp xếp các số sau theo thứ tự từ lớn đến bé: $45\,200; 45\,020; 54\,002; 45\,220$.
\choice
{$45\,020 > 45\,200 > 45\,220 > 54\,002$}
{\True $54\,002 > 45\,220 > 45\,200 > 45\,020$}
{$54\,002 > 45\,200 > 45\,220 > 45\,020$}
{$45\,220 > 54\,002 > 45\,200 > 45\,020$}
\loigiai{
Thứ tự từ lớn đến bé là: $54\,002 > 45\,220 > 45\,200 > 45\,020$.\\
\textbf{Đáp án của câu hỏi là B}
}
\end{ex}

\begin{ex} %[4D1V01]
Làm tròn số $48\,625$ đến hàng nghìn ta được
\choice
{$48\,000$}
{$48\,600$}
{$50\,000$}
{\True $49\,000$}
\loigiai{
Chữ số hàng trăm là 6 ($6 \ge 5$) nên làm tròn lên thành $49\,000$.\\
\textbf{Đáp án của câu hỏi là D}
}
\end{ex}

\begin{ex} %[4D1V01]
Tìm số tự nhiên $x$ có 5 chữ số biết chữ số hàng chục nghìn là 7 và $x$ là số chẵn lớn nhất có 5 chữ số khác nhau:
\choice
{$79\,864$}
{$78\,964$}
{\True $79\,864$}
{$79\,846$}
\loigiai{
Chữ số đầu là 7, các chữ số tiếp theo chọn lớn nhất giảm dần: nghìn là 9, trăm là 8, chục là 6, đơn vị là chẵn nên chọn 4. Số đó là $79\,864$.\\
\textbf{Đáp án của câu hỏi là C}
}
\end{ex}

\begin{ex} %[4D1V01]
Cho các chữ số $0; 2; 5; 8$. Có thể viết được bao nhiêu số có 4 chữ số khác nhau từ các chữ số đã cho?
\choice
{$24$ số}
{\True $18$ số}
{$16$ số}
{$12$ số}
\loigiai{
Hàng nghìn có 3 cách chọn (khác 0).\\
Hàng trăm có 3 cách chọn.\\
Hàng chục có 2 cách chọn.\\
Hàng đơn vị có 1 cách chọn.\\
Số lượng số viết được là: $3 \times 3 \times 2 \times 1 = 18$ số.\\
\textbf{Đáp án của câu hỏi là B}
}
\end{ex}

\begin{ex} %[4D1V01]
Số tròn chục nghìn liền trước của số $80\,000$ là
\choice
{\True $70\,000$}
{$79\,990$}
{$79\,000$}
{$90\,000$}
\loigiai{
Số tròn chục nghìn liền trước của $80\,000$ là $70\,000$.\\
\textbf{Đáp án của câu hỏi là A}
}
\end{ex}

\begin{ex} %[4D1V01]
Một thư viện có $35\,600$ cuốn sách. Sau khi nhập thêm $4400$ cuốn sách nữa thì tổng số sách của thư viện được làm tròn đến hàng chục nghìn là
\choice
{$30\,000$ cuốn}
{\True $40\,000$ cuốn}
{$50\,000$ cuốn}
{$45\,000$ cuốn}
\loigiai{
Tổng số sách: $35\,600 + 4400 = 40\,000$ cuốn.\\
\textbf{Đáp án của câu hỏi là B}
}
\end{ex}

\begin{ex} %[4D1C01]
Viết số bé nhất có 5 chữ số khác nhau mà tổng các chữ số của số đó bằng $21$.
\choice
{\True $10\,389$}
{$10\,299$}
{$12\,369$}
{$10\,479$}
\loigiai{
Để số bé nhất: chữ số hàng chục nghìn là 1, nghìn là 0, trăm là 3, chục là 8, đơn vị là 9. Tổng: $1 + 0 + 3 + 8 + 9 = 21$.\\
\textbf{Đáp án của câu hỏi là A}
}
\end{ex}

\textbf{Phần 2. Trắc nghiệm đúng sai}
\setcounterpageref{ex}{0}

\begin{ex} %[4D1H01]
Xét tính đúng sai của các khẳng định về số tự nhiên:
\choiceTF
{\True $100\,000$ là số bé nhất có 6 chữ số.}
{$99\,999$ là số lớn nhất có 6 chữ số.}
{\True Số $54\,321$ gồm $5$ chục nghìn, $4$ nghìn, $3$ trăm, $2$ chục và $1$ đơn vị.}
{\True Làm tròn số $64\,500$ đến hàng nghìn ta được $65\,000$.}
\loigiai{
\begin{itemize}
    \item $100\,000$ là số có 6 chữ số bé nhất. Đúng.
    \item $99\,999$ có 5 chữ số (không phải 6 chữ số). Sai.
    \item Cấu tạo số $54\,321$ chính xác. Đúng.
    \item Hàng trăm là 5 nên làm tròn lên $65\,000$. Đúng.
\end{itemize}
\begin{itemize}
    \item a) \textbf{Đúng}.
    \item b) \textbf{Sai}.
    \item c) \textbf{Đúng}.
    \item d) \textbf{Đúng}.
\end{itemize}
}
\end{ex}

\begin{ex} %[4D1V01]
Cho các số: $28\,500; 34\,200; 28\,050; 34\,020$. Xét tính đúng sai:
\choiceTF
{\True Số bé nhất trong dãy là $28\,050$.}
{\True Số lớn nhất trong dãy là $34\,200$.}
{Số $28\,500$ bé hơn số $28\,050$.}
{\True Sắp xếp tăng dần: $28\,050 < 28\,500 < 34\,020 < 34\,200$.}
\loigiai{
\begin{itemize}
    \item Số bé nhất: $28\,050$. Đúng.
    \item Số lớn nhất: $34\,200$. Đúng.
    \item $28\,500 > 28\,050$ nên khẳng định bé hơn là Sai.
    \item Thứ tự tăng dần chính xác. Đúng.
\end{itemize}
\begin{itemize}
    \item a) \textbf{Đúng}.
    \item b) \textbf{Đúng}.
    \item c) \textbf{Sai}.
    \item d) \textbf{Đúng}.
\end{itemize}
}
\end{ex}

\textbf{Phần 3. Trả lời ngắn}
\setcounterpageref{ex}{0}

\begin{ex} %[4D1H01]
Tìm số liền trước của số bé nhất có 5 chữ số.
\shortans{9999}
\loigiai{
\textbf{Đáp số:} $9999$\\
Số bé nhất có 5 chữ số là $10\,000$, số liền trước là $9999$.
}
\end{ex}

\begin{ex} %[4D1H01]
Viết số gồm $7$ chục nghìn, $4$ trăm và $6$ đơn vị.
\shortans{70406}
\loigiai{
\textbf{Đáp số:} $70406$\\
Số đó là $70\,406$.
}
\end{ex}

\begin{ex} %[4D1V01]
Làm tròn số $73\,480$ đến hàng trăm.
\shortans{73500}
\loigiai{
\textbf{Đáp số:} $73500$\\
Hàng chục là 8 nên làm tròn lên $73\,500$.
}
\end{ex}

\begin{ex} %[4D1V01]
Tính hiệu giữa số lớn nhất có 5 chữ số khác nhau và số bé nhất có 5 chữ số khác nhau.
\shortans{88531}
\loigiai{
\textbf{Đáp số:} $88531$\\
$98\,765 - 10\,234 = 88\,531$.
}
\end{ex}

\begin{ex} %[4D1V01]
Có bao nhiêu số tròn nghìn có 5 chữ số lớn hơn $50\,000$ và bé hơn $60\,000$?
\shortans{9}
\loigiai{
\textbf{Đáp số:} $9$\\
Các số đó là: $51\,000, 52\,000, \dots, 59\,000$ (có 9 số).
}
\end{ex}

\begin{ex} %[4D1C01]
Tìm số tự nhiên có 5 chữ số, biết rằng nếu viết thêm chữ số 2 vào bên trái số đó ta được số mới gấp 5 lần số ban đầu có thêm chữ số 2 vào bên phải.
\shortans{40816}
\loigiai{
\textbf{Đáp số:} $40816$\\
Gọi số cần tìm là $x$. Khi đó $200\,000 + x = 5 \times (10x + 2) = 50x + 10 \Rightarrow 49x = 199\,990 \Rightarrow x = 40\,816$.
}
\end{ex}
